\chapter*{Legacy Problem Audit for VI/10}
\addcontentsline{toc}{chapter}{Legacy Problem Audit for VI/10}

\noindent\textbf{Audit invariant.}
Every substantive legacy problem or exercise touching reduced/nonreduced geometry is either
included here, merged here, or explicitly routed to the canonical chapter where its main
mathematics belongs.  No problem is discarded merely to fit a target count.

\small
\begin{longtable}{p{0.14\textwidth}p{0.23\textwidth}p{0.12\textwidth}p{0.34\textwidth}}
\toprule
Legacy source & Legacy item & Disposition & Canonical handling\\
\midrule
\endfirsthead
\toprule
Legacy source & Legacy item & Disposition & Canonical handling\\
\midrule
\endhead
AG1 & Problem 1: basic opens and $D(f)=\varnothing$ iff $f$ nilpotent & MOVE & VI/08.P1; the nilpotence theorem is recalled here but the legacy problem is not duplicated.\\
AG1 & Problem 2: spectra of polynomial/arithmetic rings & MOVE & VI/07 spectrum survey.\\
AG1 & Problem 3: disconnected spectrum; idempotents; product rings & MOVE & VI/05 connectedness problem.\\
AG1 & Nilpotents are topologically invisible; $A_{\mathrm{red}}$ & INCLUDE THEORY & Core VI/10 reduction theorem and figures.\\
CA13 & Compare $k[x]/(x)$ and $k[x]/(x^2)$; same topological set, different structure & MERGE & VI/10 Examples 3--5, Exercise 16, and Problems P2/P5.\\
CA13 & Reduced versus nonreduced examples; reduced does not imply irreducible & INCLUDE THEORY & Core VI/10 comparison section and examples.\\
CA13 & Base-change warnings: prime splitting and field dependence & SPLIT & Reducedness aspect developed in VI/10.P14; systematic base change remains VI/24.\\
AG6 & Problem 6, subproblem ``a scheme which is not reduced'' & INCLUDE & VI/10.P1, preserved and expanded as a full dossier.\\
AG6 & Problem 6, remaining counterexample subproblems & MOVE & Retained for the later chapters governing irreducibility, integrality, affineness, finite type, and related properties; not duplicated here.\\
AG5 & Problem 3: scheme integral iff reduced and irreducible & MOVE & VI/27, where integral schemes and function fields are canonical.  VI/10 proves only the affine ring criterion as new P6.\\
CA volume material & Radical/nilpotent ideal exercises in commutative algebra & MOVE & Volume V/03 Radicals and Nilpotents.  VI/10 uses the results geometrically without duplicating the algebra chapter's legacy problems.\\
\bottomrule
\end{longtable}

\paragraph{Result.}
The directly relevant legacy numbered material is preserved, while mixed or broader problems
retain exactly one canonical destination.  The additional VI/10 problems are new synthesis
problems and do not replace legacy items.
